\section{The subcubic theorem}\label{sec:subcubic}
\begin{theorem}\label{thm:main}
Let $G$ be a connected simple graph with $G\not\cong K_2$ and
$\Delta(G)\le3$.  Then $\lir(\doub G)\le2$.
\end{theorem}

The proof uses finite calculations for three auxiliary state systems.
Their complete domains, recurrences and checking procedures are specified
below and in Appendix~\ref{app:cert}.  The finite calculations are not
enumerations of all graphs up to an order bound.  Inductions on paths,
trees and blocks extend the verified relations to graphs of arbitrary
order.  The only external colouring theorem used is the classical
degree-list theorem, stated as Theorem~A in \cite{CR17}.

\section{Edge weights and path interfaces}\label{sec:path}

Assign to every original edge $e$ a weight $x_e\in\{0,1,2\}$, the number
of its red copies.  Define
\begin{equation}\label{eq:actual}
 R(v)=\sum_{e\ni v}x_e,\qquad B(v)=2d_G(v)-R(v).
\end{equation}
The required condition on an edge $uv$ of weight $x$ is
\begin{equation}\label{eq:rule}
 x>0\Longrightarrow R(u)\ne R(v),\qquad
 x<2\Longrightarrow B(u)\ne B(v).
\end{equation}
When both endpoints have degree three, it is sufficient and necessary
that their red degrees differ.  The same observation applies to two
degree-two vertices, with $B=4-R$ instead of $6-R$.

Let $W=\{v:d_G(v)=3\}$ and $H=G[W]$.  Vertices of $W$ are called
\emph{roots}; a connected component of $H$ is called a \emph{core}.
Every edge outside $H$ belongs to a maximal path with internal vertices of
degree two.  Such a path is an \emph{ordinary connection} if its endpoints
are distinct roots, a \emph{return path} if both endpoints are the same
root, and a \emph{pendant path} if one endpoint is a root and the other is
a leaf.  An ordinary connection has length at least two, and a return
path has length at least three.  A pendant path can have any positive
length.  Distinct roots of an ordinary connection may belong to the same
core.  A length-one edge between roots is already an edge of $H$.

The weight of the first edge at a root is the corresponding \emph{port}.
A port is \emph{pure} if its weight is 0 or 2, and \emph{mixed} if its
weight is 1.  A root red degree is physically compatible with a port $a$
when $a\le R\le a+4$.  At a return path its two port contributions must
both be included in the same root degree.

\begin{lemma}[Path interfaces]\label{lem:path}
The following statements hold for physically compatible root degrees.
\begin{enumerate}
\item A two-edge ordinary connection with pure end ports $a,b$ is
admissible exactly when $R(u)\ne2+b$ and $R(v)\ne2+a$.
\item A three-edge ordinary connection with arbitrary ports $a,b$ is
admissible exactly when
\begin{equation}\label{eq:three}
 a\ne b,\qquad
 \neg\bigl((a=1\text{ or }b=1)\text{ and }R(u)=R(v)=3\bigr).
\end{equation}
\item Opposite pure ports are admissible on every ordinary connection of
length at least three.  Any two pure ports are admissible at length at
least four.
\item A return path accepts ports $(0,2)$ at every legal length.  If its
root red degree is nonzero modulo three, each of $(0,1)$ and $(1,2)$ is
also admissible.
\item A pendant path of length at least three accepts either pure root
port.  At length one, ports 0 and 2 exclude root red degrees 4 and 2,
respectively; at length two they exclude 2 and 4, respectively.
\end{enumerate}
\end{lemma}
\begin{proof}
For a path with weights $x_1,\ldots,x_\ell$, all comparisons between its
internal degree-two vertices are equivalent to
$x_i\ne x_{i+2}$.  At a degree-three root of red degree $r$, its incident
port $a$ permits the next weight to lie in
\begin{equation}\label{eq:P}
 P_r(a)=\{0,1,2\}\setminus
 \bigl(\{r-a:a>0\}\cup\{r-a-2:a<2\}\bigr),
\end{equation}
where values outside $\{0,1,2\}$ have no effect.
Part~1 follows by direct substitution.  For three edges, the middle
comparison is $a\ne b$ and the remaining weight must lie in
$P_{R(u)}(a)\cap P_{R(v)}(b)$.  A pure port has at least two choices.
A mixed port has at least two choices except at red degree 3, where its
set is $\{1\}$.  A pure-port set excludes 1 exactly at red degree 3.
This proves part~2.

For longer paths, separate the odd and even edge positions.  They form
two three-colour paths under the relations $x_i\ne x_{i+2}$.
At even length at least four, each position path has one fixed endpoint
and an endpoint list of size at least two.  At odd length at least five,
one position path has two fixed endpoints at distance at least two and
the other has two lists of size at least two.  Both can always be coloured.
Together with part~2 this proves part~3.  If one port is mixed but its
root red degree is not 3, its set in \eqref{eq:P} again has at least two
elements, so the same argument applies to longer paths.

For a return path set both endpoint red degrees equal to the actual
root degree.  Its nonzero residue excludes red degree 3; the preceding
argument gives part~4.  At a pure leaf port 0 or 2, the preceding weight
has allowed set $\{0,1\}$ or $\{1,2\}$, respectively.  The position-path
argument proves part~5 for length at least four.  At length three choose
the leaf port opposite to the root port and intersect two sets of size at
least two.  The length-one and length-two exclusions follow directly
from \eqref{eq:rule}.
\end{proof}

We will use the safe set
\begin{equation}\label{eq:safe}
 \safe=\{0,1,3,5,6\}.
\end{equation}
In particular, pure ports at safe roots automatically satisfy either
end of every two-edge connection.  To attach a pendant path at a modular
root of target residue $\phi$, take port 2 exactly when its length is one
and $\phi=1$, or its length is two and $\phi=2$; take port 0 otherwise.
Both choices are allowed at any length when $\phi=0$.

If $W$ is empty, $G$ is a single vertex, a path, or a cycle.
A two-edge path takes weights $(0,0)$; longer paths take first weight 0
and last weight 2 and use the position-path argument with leaf endpoint
conditions.  For a cycle, give its edge positions a proper three-colouring
in the auxiliary graph joining positions $i$ and $i+2$ modulo its length.
That auxiliary graph has maximum degree two.  The resulting weights
satisfy all comparisons.  Also, for $K_4$, on edges
$(01,02,03,12,13,23)$ take weights $(1,2,0,0,0,2)$; the red degrees are
$(3,1,4,2)$.  We henceforth exclude these cases.

\section{Three finite auxiliary lemmas}\label{sec:finite}

A core $C$ has $3-d_C(v)$ private ports at each vertex $v$.
In an auxiliary port assignment its red degree is the sum of its core
edge weights and its private port weights.  Private ports here are free
variables, not prescribed boundary weights.

\begin{lemma}[Free safe ports]\label{lem:free}
Suppose $C$ is a tree, a cactus all of whose cycle blocks are odd, or a
connected bipartite subcubic graph with at most two vertices of core
degree three.  There is an assignment for which all private ports are
pure, every red degree belongs to $\safe$, and adjacent core vertices
have distinct red degrees.
\end{lemma}

\begin{lemma}[One mixed port]\label{lem:one}
Let $C$ be a bipartite subcubic cactus.  Any specified private port can
be set to 1 so that its root red degree is 1 or 5, all other private
ports are pure, every root red degree is safe, and adjacent core vertices
have distinct red degrees.
\end{lemma}

We specify the finite calculations and their infinite extensions before
stating the third lemma.  A rooted state $(x,r)$ records the incoming
core-edge weight (or one actual free port at the final root) and the red
degree of that vertex.  Its domain is
\begin{equation}\label{eq:Omega}
 \Omega=\{(x,r):x\in\{0,1,2\},\ r\in\safe,\ x\le r\le x+4\},
 \qquad |\Omega|=9.
\end{equation}
The six rooted message sets are
\begin{align*}
 \mathcal L&=\{(0,0),(1,1),(1,3),(1,5),(2,6)\},\\
 \mathcal O&=\{(0,1),(0,3),(1,1),(1,3),(1,5),(2,3),(2,5)\},\\
 \mathcal A&=\Omega,\\
 \mathcal C&=\{(0,0),(0,1),(1,3),(2,5),(2,6)\},\\
 \mathcal D&=\{(0,0),(0,1),(0,3),(1,3),(2,3),(2,5),(2,6)\},\\
 \mathcal E&=\{(0,0),(0,1),(1,1),(1,3),(1,5),(2,5),(2,6)\}.
\end{align*}
For $c\in\{0,1,2\}$ child sets $X_1,\ldots,X_c$, define
\begin{equation}\label{eq:node}
 N(X_1,\ldots,X_c)=
 \left\{\left(x,x+\sum_{i=1}^c y_i+\sum_{j=1}^{2-c}p_j\right):
 \begin{array}{l}
 (y_i,s_i)\in X_i,\ p_j\in\{0,2\},\\
 r=x+\sum_i y_i+\sum_jp_j\in\safe,\ r\ne s_i\ (1\le i\le c)
 \end{array}\right\}.
\end{equation}
Here $x\in\{0,1,2\}$.  Every one of the $1+6+36=43$ ordered
choices has one of the six message sets as its exact result.

For a non-anchor vertex of a cycle, its third incidence is a free pure
port, denoted $F$, or one child message set $X$.  Define a Boolean matrix
$M_X$ on $\Omega$ by
\begin{equation}\label{eq:letter}
 (M_X)_{(x,r),(y,s)}=1
 \quad\Longleftrightarrow\quad
 s\ne r\text{ and }
 \begin{cases}
 s-x-y\in\{0,2\},&X=F,\\
 \exists(p,t)\in X:\ s=x+y+p,\ s\ne t,&X\ne F.
 \end{cases}
\end{equation}
The rooted interface of a cycle product $M$ consists of the states
$(z,r)$ for which there are $a,b\in\{0,1,2\}$ and $(b,s)\in\Omega$
with $r=z+a+b$, $s\ne r$, and $M_{(a,r),(b,s)}=1$.

\begin{proposition}[Finite closure identities]\label{prop:closure}
The node recurrence \eqref{eq:node} is closed on the six message sets.
Let $\mathcal P_k$ be the set of all products of $k$ letter matrices
from \eqref{eq:letter}.  Then
\begin{align*}
 (|\mathcal P_2|,|\mathcal P_4|,|\mathcal P_6|)&=(14,11,2),&
 \mathcal P_8&=\mathcal P_6,\\
 (|\mathcal P_3|,|\mathcal P_5|,|\mathcal P_7|)&=(23,6,2),&
 \mathcal P_9&=\mathcal P_7.
\end{align*}
Their rooted interface types have the following multiplicities:
\begin{center}
\begin{tabular}{ccl}
\toprule
Parity of cycle length&Number of letters&Rooted types\\
\midrule
odd&2&$\mathcal C$ once, $\mathcal D$ thirteen times\\
odd&4&$\mathcal A$ ten times, $\mathcal D$ once\\
odd&6&$\mathcal A$ twice\\
even&3&$\mathcal A$ twenty-one times, $\mathcal O$ twice\\
even&5&$\mathcal A$ six times\\
even&7&$\mathcal A$ twice\\
\bottomrule
\end{tabular}
\end{center}
Moreover the node recurrence restricted to
$\mathcal L,\mathcal O,\mathcal A$ stays within these three types.
\end{proposition}
\begin{proof}[Computer-assisted finite calculation]
Evaluate \eqref{eq:node} for the 43 choices.  Construct the seven
matrices by \eqref{eq:letter} and multiply them as Boolean relations.
The files \code{safe_cactus_certificate.py} and
\code{safe_all_cactus_certificate.py} independently recompute the
complete layers, compare sets rather than just their cardinalities,
and compute every rooted interface.  They use Python integer and set
operations only.  The complete layer sets are supplied in
\code{safe_cactus_interfaces.json} and
\code{safe_all_cactus_interfaces.json}.  Appendix~\ref{app:cert}
gives the common execution command.
\end{proof}

\begin{proof}[Proof of the tree and cactus parts of Lemmas~\ref{lem:free}
and~\ref{lem:one}]
For a tree, root at a vertex with an available private port, use that
port as the incoming incidence, and apply \eqref{eq:node} upward.
Only $\mathcal L,\mathcal O,\mathcal A$ occur.  All accept incoming
weight 0, which gives the free-port assertion, and incoming weight 1
with red degree 1 or 5, which gives the marked assertion.

For a cactus of maximum degree three, different cycle blocks cannot
share a vertex.  Contract each cycle block to a unit and induct on the
resulting tree of units.  A non-anchor cycle vertex has precisely the
third-incidence choices in \eqref{eq:letter}; two children there would
violate the degree bound.  A cycle of length $\ell$ uses $\ell-1$
letters.  The identities in Proposition~\ref{prop:closure} imply, by
induction under multiplication by any two letters, that every longer
product of the same parity has one of the listed rooted types.
Thus the argument covers arbitrary cycle lengths and numbers of blocks.

Every one of the six types accepts incoming weight 0.  This proves the
odd-cactus assertion.  In a bipartite cactus every cycle is even, and
the induction stays in $\mathcal L,\mathcal O,\mathcal A$.  Root at the
specified private port and choose $(1,1)$ or $(1,5)$ to prove
Lemma~\ref{lem:one}.  The incoming port is an actual private incidence,
not an extra edge added to the original degree.
\end{proof}

For core chains with one pure private port at each internal vertex,
write
\begin{equation}\label{eq:T}
 (x,r)\ T\ (y,s)\quad\Longleftrightarrow\quad
 r\ne s,\quad s=x+y+p\in\safe\text{ for some }p\in\{0,2\}.
\end{equation}
Both states lie in $\Omega$.  The supplied certificate gives a six-step
walk for every ordered pair of states, so
\begin{equation}\label{eq:sat}
 T^6=\Omega\times\Omega.
\end{equation}
Every state has an incoming neighbour, and therefore
$T^k=\Omega\times\Omega$ for every $k\ge6$.
Here internal chain vertices have original degree three, unlike the
degree-two internal vertices in Section~\ref{sec:path}.

\begin{proof}[Proof of the remaining part of Lemma~\ref{lem:free}]
If there are no core-degree-three vertices, the core is a single vertex,
a path, or an even cycle, already covered above.  Otherwise compress its
core-degree-two chains.  If $n_1,n_3$ count core degrees one and three,
and $\beta=|E(C)|-|V(C)|+1$, then
\begin{equation}\label{eq:euler}
 n_1=n_3+2-2\beta.
\end{equation}
With one degree-three vertex the only shapes are a three-arm tree and
an even cycle with an arm.  With two, the only shapes are a double-star
tree, an even cycle attached to a fork with two arms, an even cycle with
an arm at each of its two forks, a theta, or two even cycles joined by
a chain.  These follow by counting loops and paths in the compressed
shape, whose two forks have degree three.

There are 163 complete finite shape and length representatives, with
their domain specified in Table~\ref{tab:sparse}.  Every representative
has a witness assigning odd safe red degree to each fork, consistent
with its three incident weights.  The checker
\code{sparse_bip_core_certificate.py} verifies the complete domain,
all fork sums and all chain relations, as well as \eqref{eq:sat}.

Keep a length $n\le6$, replace a longer odd length by 7, and a longer
even length by 8.  Any chain of length at least seven has at least six
internal transitions, so \eqref{eq:sat} lifts its witness with unchanged
end weights and fork degrees.  Parities and simplicity restrictions
are preserved.  An arm can have arbitrary positive length: an odd-red
fork can be completed immediately at a leaf by two pure ports; from
every state one step reaches a state that admits such a completion,
and every state has an outgoing step.  These finite assertions are also
checked.  Thus all arm lengths and all chain lengths are covered.
\end{proof}

\begin{definition}[Four-port tree]\label{def:four}
A four-port tree is a subcubic tree core with two distinguished vertices
$a,b$ at odd distance.  Neither has a short port.  On each side of the
$a$--$b$ path there are exactly two positions, each either a direct long
private port at the distinguished vertex or an arm ending at a leaf
with one long and one short port.  Every other internal vertex of such
an arm or of the $a$--$b$ path has one short private port.  Thus there
are precisely four long ports.  The distinguished vertices have no
other incidences.
\end{definition}

Figure~\ref{fig:four} shows one possible topology.  A position at a
distinguished vertex may instead be a direct long port.
\begin{figure}[htbp]
\centering
\setlength{\unitlength}{0.9mm}
\begin{picture}(140,54)
\put(35,25){\circle*{2.4}}\put(105,25){\circle*{2.4}}
\put(35,25){\line(1,0){70}}
\put(55,25){\circle*{1.5}}\put(85,25){\circle*{1.5}}
\put(55,25){\line(0,1){8}}\put(85,25){\line(0,1){8}}
\put(55,36){\makebox(0,0){$s$}}\put(85,36){\makebox(0,0){$s$}}
\put(35,25){\line(-4,3){20}}\put(35,25){\line(-4,-3){20}}
\put(105,25){\line(4,3){20}}\put(105,25){\line(4,-3){20}}
\put(15,40){\circle*{1.5}}\put(15,10){\circle*{1.5}}
\put(125,40){\circle*{1.5}}\put(125,10){\circle*{1.5}}
\put(15,40){\line(-1,0){13}}\put(15,10){\line(-1,0){13}}
\put(125,40){\line(1,0){13}}\put(125,10){\line(1,0){13}}
\put(15,40){\line(0,1){8}}\put(15,10){\line(0,-1){8}}
\put(125,40){\line(0,1){8}}\put(125,10){\line(0,-1){8}}
\put(15,51){\makebox(0,0){$s$}}\put(15,0){\makebox(0,0){$s$}}
\put(125,51){\makebox(0,0){$s$}}\put(125,0){\makebox(0,0){$s$}}
\put(2,44){\makebox(0,0){$\ell_1$}}\put(2,6){\makebox(0,0){$\ell_2$}}
\put(138,44){\makebox(0,0){$\ell_3$}}\put(138,6){\makebox(0,0){$\ell_4$}}
\put(35,18){\makebox(0,0){$a$}}\put(105,18){\makebox(0,0){$b$}}
\end{picture}
\caption{A four-port tree with two distinguished vertices $a,b$, four
long ports $\ell_i$, and short ports labelled $s$.  The middle path
is odd; the arms and middle path may be longer than shown, with one
short port at every added internal core vertex.  All vertices drawn
have original degree three.  Only $a,b$ have no short incidence and
may use red degrees 2 or 4.}
\label{fig:four}
\end{figure}

\begin{lemma}[Two mixed ports]\label{lem:two}
In any four-port tree, any two distinct long ports can be set to 1 so
that their roots have red degree different from 3.  All other private
ports can be pure, adjacent core vertices can have distinct red
degrees, and every vertex carrying a short port can have safe red
degree.  The two distinguished vertices may have red degree 2 or 4.
\end{lemma}
\begin{proof}
For an arm, its incoming message at the leaf is $\mathcal L$ if the
long port is pure, and
$Q=\{(0,1),(2,5)\}$ if the long port is marked.  At an internal
vertex the message map is
\begin{equation}\label{eq:U}
 U(X)=\{(x,x+y+p):(y,s)\in X,\ p\in\{0,2\},\
 x+y+p\in\safe,\ x+y+p\ne s\}.
\end{equation}
Exact iteration gives cardinalities $5,7,9,9$ from $\mathcal L$, and
$2,3,5,7,7,9,9$ from $Q$.  In each case the last two sets are equal
to $\Omega$, not merely equinumerous.  Thus an unmarked position has
four types, direct or arm lengths $1,2,\ge3$; a marked position has
seven types, direct or arm lengths $1,2,3,4,5,\ge6$.

Up to exchanging the two ends, the mark distribution is $(2,0)$ or
$(1,1)$.  The first has
$\binom82\binom52=280$ choices of unordered pairs of position types;
the second has $(7\cdot4)^2=784$ choices.  Two direct marked ports at
the same distinguished vertex are included.  The connecting path is
odd, and its representatives have lengths $1,3,5,7$.  Hence the complete
domain contains $(280+784)\cdot4=4256$ cases.

For each case, \code{weak_two_marked_interfaces.json} gives the end
states, the two position choices at each end, and compatibility with
the connecting path.  The checker recomputes the domain, arm sets,
actual sums and connecting relations; every case passes.  A second
verifier recomputes existence from these definitions without using the
stored witnesses.  Details appear in Appendix~\ref{app:cert}.

Arm fixed points prove lifting at every larger arm length.  For the
middle path, \eqref{eq:sat} proves lifting at every length at least
seven if the two end degrees are safe.  If an end red degree is 2 or 4,
it differs from every internal safe red degree; a word valid with the
same end-edge weight and virtual end red degree 3 is still valid with
the actual degree.  The virtual state is physical for every possible
end weight.  Thus larger odd lengths preserve both actual end degrees
and end weights.  These observations prove the lemma for arbitrary
four-port trees, rather than only for the finite representatives.
\end{proof}

\section{Classifying cores and counting long ports}\label{sec:class}

Call an ordinary connection \emph{short} if its length is two, and
\emph{long} if its length is at least three.  For a bipartite core $C$
with parts $A,B$, count the following objects on each part:
\begin{enumerate}
\item every pendant port rooted in that part;
\item every endpoint in that part of a short connection whose two roots
both belong to $C$;
\item every vertex in that part incident with neither a pendant path
nor any short connection, including short connections to other cores.
\end{enumerate}
Let the resulting counts be $k_A,k_B$.  In item~3 a core-degree-three
vertex is always counted, even though it has no private port.

Assign cores to classes $\SC$ and $\MC$ by the following fixed rules,
before choosing any weights.  Put in $\MC$ every core with a return
path, every nonbipartite core incident with an ordinary three-edge
connection, and every bipartite core incident with an ordinary
three-edge connection and satisfying $\max(k_A,k_B)\ge2$.
Among the remaining cores, put in $\SC$ those covered by
Lemma~\ref{lem:free}, and put the rest in $\MC$.

The letters specify how a core will be handled, not that all degrees
in its final colouring must have a common residue.  In particular,
Lemma~\ref{lem:two} may give red degree 2 or 4 at vertices of an
$\SC$ core that carry no short port.

\begin{lemma}\label{lem:count}
An $\SC$ core incident with an ordinary three-edge connection is
bipartite and has at most four long ports.  If it has four ports that
belong to three-edge connections between $\SC$ cores, it is a four-port
tree as in Definition~\ref{def:four}.
\end{lemma}
\begin{proof}
Nonbipartite cores incident with such a connection have been assigned
to $\MC$.  For a remaining bipartite core, $k_A,k_B\le1$.
Let $p$ be the number of pendant ports, $s$ the number of short
connections internal to $C$, and $t$ the number of vertices in item~3.
Then
\begin{equation}\label{eq:budget}
 t+p+2s=k_A+k_B\le2.
\end{equation}
There is no return path.  A core leaf has at most one long port, plus
one more if it is counted in $t$.  A core-degree-two vertex with a long
port is counted in $t$.  A core-degree-three vertex is counted in $t$
but has no long port.  Therefore, writing $L$ for the total number of
long ports, \eqref{eq:euler} gives, for a nontrivial core,
\begin{equation}\label{eq:Lbound}
 L\le n_1+t-n_3=2-2\beta+t
 \le4-2\beta-p-2s\le4.
\end{equation}
A one-vertex core has only three ports.

If four of these ports belong to the stated three-edge connections,
all inequalities are equalities.  Thus $\beta=p=s=0$, $t=2$, and the
two vertices counted in $t$ lie in opposite parts.  All core-degree-three
vertices are among them.  Their connecting path has odd length.  Each
other internal core vertex has its one short port; each other core
leaf has one long and one short port.  Removing the middle path leaves
exactly two positions at each end, either direct ports or unbranched
arms.  There are no other long ports, pendant paths or internal short
connections.  This is precisely Definition~\ref{def:four}.
\end{proof}

\begin{corollary}\label{cor:cactus}
Every $\SC$ core incident with an ordinary three-edge connection is a
bipartite cactus.
\end{corollary}
\begin{proof}
Its $L$ is positive, so \eqref{eq:Lbound} gives $\beta\le1$.
A connected graph with at most one cycle is a cactus.
\end{proof}

\section{Making all safe three-edge connections compatible}\label{sec:glue}

Form an indexed multigraph $F$ whose vertices are the $\SC$ cores.
Every ordinary three-edge connection with both ends in $\SC$ gives an
edge of $F$.  Roots in the same core give a loop, and different
connections may give parallel edges.  Loops contribute two to degree.
By Lemma~\ref{lem:count}, $\Delta(F)\le4$.

\begin{lemma}\label{lem:orient}
There is a spanning forest $T$ of $F$ and an orientation of
$D=E(F)\setminus E(T)$ with outdegree at most two at each vertex.
Outdegree two is possible only at a vertex of $F$ of degree four.
\end{lemma}
\begin{proof}
Take any spanning forest.  Orient $D$ so that its outdegree at each
vertex lies between the floor and ceiling of half its degree.  Such
an orientation is obtained by joining each odd-degree vertex to a
temporary vertex, orienting Euler tours in the resulting even-degree
multigraph, and deleting the temporary edges.  A loop has one tail.

In a component of $F$ with at least two vertices, every vertex has an
incident forest edge.  Thus its $D$-degree is at most three, and at
most two if its $F$-degree is at most three.  In a singleton component
all edges are loops, and at most two loops are possible.  These facts
prove both assertions.
\end{proof}

For every $\SC$ core of outdegree zero use Lemma~\ref{lem:free}.
For outdegree one, mark its outgoing half-port by 1 and use
Corollary~\ref{cor:cactus} and Lemma~\ref{lem:one}.
For outdegree two, Lemma~\ref{lem:count} gives a four-port tree, to
which Lemma~\ref{lem:two} applies.  Only long ports are marked.
Every short-incident root remains safe.

Each $D$ edge now has exactly one mixed and one pure port, with red
degree different from 3 at its mixed end.  It is compatible by
Lemma~\ref{lem:path}.  A loop in $F$ uses two distinct actual roots of
the same core; it is not an original return path.

Every $T$ edge still has pure ports $a,b$.  Choose a colour-complement
bit $\eta_C$ for each core.  Complementation changes every core weight
and private port from $x$ to $2-x$ and changes each root red degree
from $r$ to $6-r$.  To make the two ports of a $T$ edge opposite impose
\begin{equation}\label{eq:xor}
 \eta_C\oplus\eta_{C'}=1\oplus(a/2)\oplus(b/2).
\end{equation}
These equations are simultaneously solvable on a forest.
Complementation preserves safety at short ports, unequal red degrees
on core edges, and the condition that a marked root is not of red
degree 3.  Thus all $D$ connections remain compatible.

All ordinary three-edge connections between $\SC$ cores are now
compatible.  At an $\SC$--$\MC$ three-edge connection the $\SC$ port
is pure, since marking was confined to $D$.  Set the $\MC$ port
opposite to it.  At a three-edge connection between two $\MC$ roots
use ports 0 and 2.  At longer ordinary connections any $\MC$ port
can be set to 0.  These choices put no restriction on the eventual
physical red degrees of $\MC$ roots.

\section{Targets and realizability modulo three}\label{sec:mod}

All arithmetic on target values in this section is in $\F$.

\begin{lemma}[Realizing sums]\label{lem:linear}
Let $C$ be a connected simple graph and $q:V(C)\to\F$.
There are weights in $\{0,1,2\}$ whose incident sums equal $q$ modulo
three if and only if, in the bipartite case with parts $A,B$,
\begin{equation}\label{eq:balance}
 \sum_Aq-\sum_Bq=0.
\end{equation}
There is no constraint if $C$ is nonbipartite.
\end{lemma}
\begin{proof}
In a bipartite graph, each edge contributes once to each part, proving
necessity.  Fix a spanning tree.  Setting non-tree weights to zero,
remove leaves successively and set their parent edge to the remaining
target.  The final equation is exactly \eqref{eq:balance}.
For a nonbipartite graph, sign vertices alternately along the tree.
Some non-tree edge has equal signs at its ends.  Its coefficient in
the signed defect is $2$ or $-2$, invertible in $\F$.  Choose its
weight to cancel the defect and then eliminate along the tree.
Choose representatives 0,1,2 for the resulting field elements.
\end{proof}

For an $\MC$ core, let $c(v)$ be the sum of its private ports and seek
a proper target colouring $\phi(v)\in\F$ of its core edges, then
realize the core sums $q(v)=\phi(v)-c(v)$.  The final actual degree
satisfies $R(v)\equiv\phi(v)$, so adjacent core degrees differ.

If a short connection joins $v$ to an $\SC$ root with port $a$, delete
$(2+a)\bmod3$ from the list at $v$.  Denote the resulting list by
$L(v)$.  Since $a$ is pure, only colours 1 or 2 are deleted; hence
\begin{equation}\label{eq:lists}
 0\in L(v),\qquad |L(v)|\ge d_C(v).
\end{equation}
Indeed at most $3-d_C(v)$ private ports can impose deletions.
A return path or a three-edge port at a vertex leaves a strict
surplus $|L(v)|>d_C(v)$ there.

We use the following established theorem \cite[Theorem~A]{CR17}:
a connected graph that is not a Gallai tree is colourable from any
lists of sizes at least the vertex degrees.  A Gallai tree is a
connected graph whose blocks are complete graphs or odd cycles.
Also, a single strict-surplus vertex suffices in any connected graph:
root a spanning tree there and colour in reverse tree order.

\subsection{Nonbipartite cores}

If a nonbipartite $\MC$ core has a return path or an ordinary
three-edge connection, it has a strict-surplus vertex.  Otherwise,
if it were a Gallai tree, its blocks would be bridges, odd cycles,
or $K_4$.  A $K_4$ uses all degree-three incidences, forcing
$G=K_4$, already treated.  Without it the core is an odd-cycle
cactus and would have been assigned to $\SC$.  Thus the remaining
core is not a Gallai tree.  In all cases a proper list target exists,
and Lemma~\ref{lem:linear} realizes it for any subsequent private
offsets.  No bipartite defect needs to be balanced.

\subsection{Bipartite cores with at least two safe short exits}

Let $f$ count short connections from $C$ to $\SC$ cores, with
multiplicity.  If $f\ge2$, choose a proper target from the lists.
This is justified as follows.  If $C$ contains a cycle, its bipartite
structure makes it a non-Gallai graph.  If it is a tree, its assignment
to $\MC$ implies the presence of a return path or a three-edge
connection, giving a strict-surplus vertex.

Reserve this core's own port on each of two safe short exits.  Each
can be set independently to 0 or 2: the opposite root is safe and
therefore accepts either value.  Each changes the defect in
\eqref{eq:balance} by a nonzero element of $\F$, without changing
any target.  Two such binary changes cover all residues and will
balance the defect after other offsets have been fixed.

\subsection{Bipartite cores with at most one safe short exit}

Suppose $f\le1$.  If $C$ has a return path, choose its root in part
$B$, set $\phi=0$ on $A$ and $\phi=t$ on $B$, and choose
$t\in\{1,2\}$ to respect the possible single list restriction.
The return path accepts all three port pairs $(0,1),(0,2),(1,2)$,
because its root target is nonzero.  Their total contributions
$1,2,3$ cover the three defects.

Now suppose there is no return path.  If $C$ was assigned to $\MC$
by the $k$-count rule, some side has at least two objects counted in
$k$.  Any other bipartite $\MC$ core not already covered has at
least three core-degree-three vertices; otherwise
Lemma~\ref{lem:free} would have placed it in $\SC$.  One side
therefore has two such vertices, each counted by item~3.  Again one
side has at least two counted objects.

Take this side as $A$, put $\phi=0$ on $A$, and put
$\phi=t\in\{1,2\}$ on $B$.  Since all lists contain 0 and there
is at most one deletion in total, a choice of $t$ respects all lists.
Two objects on $A$ will provide independent defect adjustments.

\section{Balancing all cores without disturbing short connections}
\label{sec:controls}

First fix the target at every root incident with a short connection.
For a short connection between two $\MC$ roots, give either endpoint
the port
\begin{equation}\label{eq:shortchoice}
 p=\begin{cases}2,&\text{the other root has target }2,\\
                  0,&\text{otherwise}.
   \end{cases}
\end{equation}
Then the other root target differs from $(2+p)\bmod3$, as required
by Lemma~\ref{lem:path}.  Give each $\MC$ pendant its port from
Section~\ref{sec:path}; give each return path the default pair
$(0,2)$ until a selected return adjustment is made.

Consider the remaining case from Section~\ref{sec:mod}: a bipartite
core with targets 0 on $A$, $t$ on $B$ and two counted objects on
$A$.  Define
\begin{equation}\label{eq:defect}
 \delta_C=\sum_{v\in A}(\phi(v)-c(v))-
           \sum_{v\in B}(\phi(v)-c(v)).
\end{equation}
Its counted objects have the following interpretations.
\begin{enumerate}
\item A pendant port at a zero-target root may take either pure
value, at every length, by Lemma~\ref{lem:path}.
\item For an internal short connection whose endpoint $u\in A$
has target 0, the \emph{other endpoint's own port} is free to take
0 or 2.  The actual red degree at $u$ is 0 modulo three, hence
avoids both 2 and 4.  The comparison at the other endpoint is
protected by the already chosen port at $u$ in \eqref{eq:shortchoice}.
\item A zero-target root with no pendant and no short incidence may
change its target from 0 to $3-t$.  This remains different from
all neighbouring core targets $t$.  No short port or pendant choice
depends on this target.
\end{enumerate}

Each object changes \eqref{eq:defect} by a nonzero element of $\F$.
The variable freed by item~2 is at the opposite endpoint, so its
sign in the defect can be either sign; nonzeroness is all that is
needed.  Different short endpoints free different actual ports.
If both ends are on $A$, both actual ports are independently free.
An item~3 vertex has no short or pendant incidence, so it cannot
overlap variables of the first two kinds.  Thus two counted objects
really give independent binary choices.

For any nonzero $\alpha,\beta\in\F$,
\begin{equation}\label{eq:binary}
 \{\alpha\epsilon+\beta\epsilon':
          \epsilon,\epsilon'\in\{0,1\}\}=\F.
\end{equation}
Consequently the two chosen objects can cancel any initial defect.
The two-safe-exit case uses the same calculation.  The return case
uses its three total contributions instead.

All cores can be balanced separately.  Changing a target in item~3
never changes a short endpoint or pendant target.  Changing a port
uses only a pendant, a short connection internal to the same core,
or this core's end of a connection to an $\SC$ root.  It never
changes another $\MC$ core's offsets.  On a safe short exit the
opposite root accepts either pure port, and its own degree does not
change.  Long path interfaces accept arbitrary actual root degrees.
There is therefore no remaining coupled system between different
modular cores.

Apply Lemma~\ref{lem:linear} to each core with its balanced targets
and final offsets.  All resulting core weights belong to $\{0,1,2\}$.
The actual root degrees are the sums of these incident core weights
and the fixed ports, so they are physical and satisfy
$R(v)\equiv\phi(v)$.  Every direct core edge has distinct red
degrees at its ends and hence distinct blue degrees as well.

\section{Completion of the proof}\label{sec:complete}

\begin{proof}[Proof of Theorem~\ref{thm:main}]
The cases with no degree-three root and the case $K_4$ were treated
in Section~\ref{sec:path}.  Otherwise decompose the graph into cores
and paths and classify its cores as in Section~\ref{sec:class}.
Sections~\ref{sec:glue}--\ref{sec:controls} assign all core weights
and all root ports, with unequal degrees along every direct core
edge.  It remains to fill the path interiors.

For a short connection between two $\SC$ roots, both roots carrying
short ports are safe.  Any exceptional red degree 2 or 4 allowed by
Lemma~\ref{lem:two} occurs only at a root without short ports.
For an $\SC$--$\MC$ short connection, the safe root's comparison
is automatic and the modular root's comparison follows from its list.
For an $\MC$--$\MC$ short connection, \eqref{eq:shortchoice} and
the frozen short targets ensure both comparisons; Section~\ref{sec:controls}
preserves them.

Every ordinary three-edge connection either has opposite pure ports
or is a $D$ edge with one mixed port whose root degree is not 3.
Every longer ordinary connection has pure ports.  Return and pendant
paths have their admissible pairs or ports.  Lemma~\ref{lem:path}
therefore fills every path.

The interiors are edge-disjoint, and their first and last weights
were already counted in the root degrees.  Filling them does not
alter a root degree.  This verifies core edges, internal path edges,
root-end edges and leaf-end edges, exhausting $E(G)$.  Replacing a
weight $x$ by $x$ red and $2-x$ blue copies yields \eqref{eq:rule}
on every edge of the doubled graph.
\end{proof}

The argument uses the common identity $B=6-R$ at the roots, their
three-incidence budget and the resulting bound $\Delta(F)\le4$.
No theorem for graphs of larger maximum degree is asserted here.

